\label{chap:83-03}

\section{Genealogical construction of the continuum}

\subsection{APP, TRIT and TPK: governing definitions}

The construction begins with three irreducible and causally ordered objects. The object of \emph{All Possible Paths} (hereafter, \APP) constructs the common nonadic support, its graph of all finite paths, and the two sum--product evaluation sheets. The tritic unit (hereafter, \TRIT) is the minimal local unit of topological, geometric, informational, and, once ordered, temporal relation: its visible signature has three states, and its lift preserves the carry that a binary or residual reading would lose. The \emph{Topological Phase Kernel} (hereafter, \TPK) is the transport category that composes these units into routes and preserves, within the same state, sheet, orientation, boundary, genealogical ledger, memory, phase, and survival.

\begin{center}
\begin{tabularx}{.94\textwidth}{@{}p{.10\textwidth}Y Y@{}}
\toprule
Object & Construction & Invariantes transportados\\
\midrule
APP & positive chart, finite torus, Cayley graph, paths, and leaves
\(\Sigma,\Pi\) & marker, cell, neighborhood, sum, product, quotient and closure\\
TRIT & ternary signature visible and local topological-temporal lift &
residue, carry, sheet, orientation, threshold and reversal\\
TPK & category of paths, fibers, and compatible extensions & cylinder,  
boundary, genealogical record, memory, phase, holonomy and survival\\
\bottomrule
\end{tabularx}
\end{center}

The exact construction of APP is established on a marked segment and a
finite chart; the real line appears afterward as a limiting publication. The
reference interval, its marks, and the oriented path that traverses them are:
\begin{equation}
I_{10}=[0,10],\qquad
\mathcal M_{10}=I_{10}\cap\mathbb Z=\{0,1,\ldots,10\},
\qquad
\mathsf P_{10}:0\longrightarrow1\longrightarrow\cdots
\longrightarrow9\longrightarrow10,
\label{p12-83-c3-s1:eq:app-marked-segment}
\end{equation}
The endpoints fix the edge and the nine interior marks
\(I_{10}^{\circ}\cap\mathbb Z\) form the positive alphabet. The exact
mathematical object of departure is the chart
\begin{equation}
\mathcal D_9=\{1,\ldots,9\},
\qquad \mathcal D_8=\{1,\ldots,8\},
\qquad
\lambda_9:\mathcal D_9\xrightarrow{\sim}\mathbb Z/9\mathbb Z,
\qquad \lambda_9(9)=[0]_9.
\label{p12-83-c3-s1:eq:app-positive-chart}
\end{equation}
The mark \(9\) is simultaneously the last positive representative and the
publication of the null class. For every \(n\geq1\), the positive digital root and
quotient are separated by means of
\begin{equation}
\operatorname{dr}_9(n)=1+((n-1)\bmod9),
\qquad q_9(n)=\left\lfloor\frac{n-1}{9}\right\rfloor,
\qquad n=9q_9(n)+\operatorname{dr}_9(n).
\label{p12-83-c3-s1:eq:positive-digital-root-carry}
\end{equation}
To the path \(1\to2\to\cdots\to9\) we add the unique closing edge \(9\to1\). The successor and its memory are separated:
\begin{equation}
u_+(a)=\begin{cases}a+1,&a<9,\\1,&a=9,\end{cases}
\qquad
\kappa_+(a)=\begin{cases}0,&a<9,\\1,&a=9,\end{cases}
\qquad a+1=u_+(a)+9\kappa_+(a).
\label{p12-83-c3-s1:eq:app-successor-carry}
\end{equation}
For that reason \(9\to1\) closes the phase, increases the memory, and prolongs the
genealogical state.

The square of the chart produces the finite torus and its movement graph:
\begin{equation}
X_9=(\mathbb Z/9\mathbb Z)^2,
\qquad
\Gamma_{\rm APP}=\operatorname{Cay}
\bigl(X_9,\{\pm e_1,\pm e_2\}\bigr),
\qquad
\operatorname{Path}(\Gamma_{\rm APP}).
\label{p12-83-c3-s1:eq:app-cayley-torus}
\end{equation}
Identifying opposite sides preserves the cardinality of the chart and adds the continuation law. The core of \(8^2=64\) interstices is completed by the minimal gnomon that publishes the zero class in both coordinates:
\begin{equation}
\mathcal D_9^2=\mathcal D_8^2\sqcup\mathcal G_9,
\qquad
81=64+(9+8)=64+17.
\label{p12-83-c3-s1:eq:app-gnomon-public}
\end{equation}
The same law is prolonged to charts of greater depth; the scale changes,
but the closure, the shared edge, and the memory of the turn preserve their type.

The associated volumetric completion is expressed through the strata identity.
\begin{equation}
10^3=(9+1)^3=729+243+27+1.
\label{p12-83-c3-s1:eq:completion-unit}
\end{equation}
\input{figures/fig_app_emrh_mini}
The interior volume, faces, edges, and point closure are thus
typed as distinct layers of the same completed unit. The evaluations act on the 81
cells of the positive chart
\[
\Sigma,\Pi:\mathcal D_9^2\longrightarrow\mathcal D_9,
\qquad
\Sigma(i,j)=\operatorname{dr}_9(i+j),
\qquad
\Pi(i,j)=\operatorname{dr}_9(ij),
\]
Their residues and quotients raised are defined by
\[
\begin{aligned}
R^+(i,j)&:=\Sigma(i,j),&
Q^+(i,j)&:=\frac{i+j-R^+(i,j)}9,\\
R^\times(i,j)&:=\Pi(i,j),&
Q^\times(i,j)&:=\frac{ij-R^\times(i,j)}9,
\end{aligned}
\]
and, therefore,
\[
\begin{aligned}
\widehat\Sigma(i,j)&:=\bigl(R^+(i,j),Q^+(i,j)\bigr),\\
\widehat\Pi(i,j)&:=\bigl(R^\times(i,j),Q^\times(i,j)\bigr),
\end{aligned}
\qquad
\widehat\Sigma,\widehat\Pi:
\mathcal D_9^2\longrightarrow\mathcal D_9\times\mathbb Z_{\geq0}.
\]
The bijection \(\lambda_9^{\times2}:\mathcal D_9^2\to X_9\) transports
these operations and their lifts to the torus. Therefore, APP is the tuple
\begin{equation}
\APP=\bigl(\mathcal D_9,\lambda_9,X_9,\Gamma_{\rm APP},
\operatorname{Path}(\Gamma_{\rm APP}),
\Sigma,\Pi,\widehat\Sigma,\widehat\Pi\bigr),
\label{p12-83-c3-s1:eq:app-full-definition-public}
\end{equation}
The table constitutes a two-dimensional representation of this complete tuple.
The additive sheet is Latin; the multiplicative sheet conserves
the ternary radical of the ring \(\mathbb Z/9\mathbb Z\)
\begin{equation}
R_9:=\mathbb Z/9\mathbb Z,
\qquad
\mathfrak m:=([3]_9)=\{[0]_9,[3]_9,[6]_9\}\triangleleft R_9,
\qquad \mathfrak m^2=(0).
\label{p12-83-c3-s1:eq:app-radical-369}
\end{equation}
The internal rotation is the automorphism
\[
\rho:\mathfrak m^3\longrightarrow\mathfrak m^3,
\qquad \rho(x_1,x_2,x_3)=(x_2,x_3,x_1),
\]
whose distinguished orbit is
\[
([3]_9,[6]_9,[0]_9)\xrightarrow{\rho}
([6]_9,[0]_9,[3]_9)\xrightarrow{\rho}
([0]_9,[3]_9,[6]_9)\xrightarrow{\rho}
([3]_9,[6]_9,[0]_9).
\]
The direct reduction sends \(3,6,9\) to \(0,0,0\). The internal coordinate
of the ideal, by contrast, conserves the extracted factor:
\begin{equation}
\eta_{\mathfrak m}:(\mathfrak m,+)
\xrightarrow{\;\sim\;}(\mathbb Z/3\mathbb Z,+),
\qquad
\eta_{\mathfrak m}([3a]_9):=[a]_3,
\qquad
([3]_9,[6]_9,[0]_9)
\xmapsto{\;\eta_{\mathfrak m}^{\times3}\;}
([1]_3,[2]_3,[0]_3),
\label{p12-83-c3-s1:eq:app-369-120}
\end{equation}
so that the \(3,6,9\) sector remains intact and its internal coordinate is
read as \(1,2,0\). This map belongs to the radical of APP; the
temporal TRIT projection constitutes the subsequent calendrical reader. The
three calendrical classes
remain separated as
\[
\mathcal C_+=\{1,4,7\},\qquad
\mathcal C_-=\{2,5,8\},\qquad
\mathcal C_0=\{3,6,9\}.
\]
The incompatibility between the two leaves becomes a measurable invariant. On
\(\mathscr A_3=\mathcal D_9^3\), with uniform measure
\(\nu_3(x)=9^{-3}\), let
\(\mathcal H_3=L^2(\mathscr A_3,\nu_3)\). We define
\[
\Sigma_3,\Pi_3:\mathscr A_3\longrightarrow\mathcal D_9,
\qquad
\Sigma_3(x)=\operatorname{dr}_9(x_1+x_2+x_3),
\qquad
\Pi_3(x)=\operatorname{dr}_9(x_1x_2x_3),
\]
and the orthogonal projectors
\[
P_{\Sigma,3}:=\mathbb E_{\nu_3}[\,\cdot\mid\sigma(\Sigma_3)],
\qquad
P_{\Pi,3}:=\mathbb E_{\nu_3}[\,\cdot\mid\sigma(\Pi_3)]
\in\mathcal B(\mathcal H_3).
\]
The joint enumeration of the \(9^3\) words then determines
\[
\bigl\|[P_{\Sigma,3},P_{\Pi,3}]\bigr\|=\frac{\sqrt3}{4},
\qquad
B_c=7I+2R=9P_++5P_-,
\qquad
459-405=54.
\]
The first number reappears as the electronic prefactor; the second
organizes the central block; the third records the global sum--product
defect.

\subsection{TRIT: local topological unit and temporal relation unit}

HMT first fixes the visible signature
\begin{equation}
\mathrm{TRIT}_{\rm vis}=\{-1,0,+1\},
\qquad \tau\in\mathrm{TRIT}_{\rm vis}.
\label{p12-83-c3-s1:eq:trit-visible-signature}
\end{equation}
In contrast to the binary unit, the TRIT preserves three local regimes. In
the uniform distribution, its information and entropy are distinguished as
\begin{equation}
I_{\rm TRIT}=\log3\ \text{nats}=\log_2 3\ \text{bits},
\qquad S_{\rm TRIT}=k_BI_{\rm TRIT}.
\label{p12-83-c3-s1:eq:trit-information-unit}
\end{equation}
The bit remains the binary unit; the TRIT is the minimal unit capable of simultaneously distinguishing return, threshold, and opening. The equality \eqref{p12-83-c3-s1:eq:trit-information-unit} quantifies the information of the triad; the cost \(k_B\Theta\log3\) belongs to the subsequent thermal reset reader. Dimensional Mechanics also preserves the arithmetic antecedent of that signature. For \((i,j)\in\mathcal D_9^2\), the associated lifted arithmetic cell is
\begin{equation}
c_{\rm MD}(i,j;\tau)=
\bigl(i,j;R^+,Q^+,R^\times,Q^\times;\tau\bigr).
\label{p12-83-c3-s1:eq:trit-complete-local-unit}
\end{equation}
The residues \(R^+,R^\times\) are the visible publication of the two sheets;
the quotients \(Q^+,Q^\times\) conserve the exact memory of Euclidean
division; and \(\tau\) orients the transition. The visible signature occupies the first
level; the cell that bears it adds the arithmetic antecedents, and the
enriched state then completes boundary, route, phase, and survival.

It is useful to distinguish four strata. The visible signature is \(\tau\); the
residue--carry pair is its local lift; \(c_{\rm MD}\) is the APP
arithmetic cell equipped with that signature; and the TPK state adds extension,
boundary, genealogical ledger, phase, and survival. This hierarchy preserves the distinct
identity of the observed coordinate, its lift, the carrier cell, and the
enriched history that transports it.

In the balanced ternary chart, the lifted local integer
\(n\in\{-4,-3,\ldots,3,4\}\) decomposes uniquely as
\begin{equation}
\widehat n=(r,\kappa),
\qquad n=r+3\kappa,
\qquad r,\kappa\in\{-1,0,+1\}.
\label{p12-83-c3-s1:eq:trit-balanced-lift}
\end{equation}
The visible signature is \(\tau=r\); \(\kappa\) conserves the carry that the projection
\(n\mapsto r\) would forget.
The fiber over the visible zero contains three distinct states,
\[
0^+=(0,+1),\qquad 0^0=(0,0),\qquad 0^-=(0,-1),
\]
because \(n=-3,0,3\) have the same visible residue and different carries. This
is the minimal topological lift of MD: an oriented APP transition
is endowed with the information required to invert it and to decide which
prolongations remain compatible.

The signature also has a sheet-memory rule. If \(m_k\) is the active
arithmetic mode, then
\begin{equation}
m_{k+1}=
\begin{cases}
\Sigma,&\tau_k=+1,\\
\Pi,&\tau_k=-1,\\
m_k,&\tau_k=0.
\end{cases}
\label{p12-83-c3-s1:eq:trit-zero-memory}
\end{equation}
Zero encodes retention of the previous sheet. For that reason the word
\((+1,-1,0)\) publishes the sequence \((\Sigma,\Pi,\Pi)\): the third state
conserves the active multiplicative mode.

The temporal relation is obtained by orderly parametrizing these cells.
Each signature carries a quadratic algebra.
\begin{equation}
\mathbb A_\tau=\mathbb R[X]/(X^2+\tau),
\qquad J_\tau=[X],
\label{p12-83-c3-s1:eq:trit-quadratic-algebra-local}
\end{equation}
so that
\[
\begin{array}{c@{\qquad}c@{\qquad}l}
\toprule
\tau & J_\tau^2 & \text{ordered temporal relation}\\
\midrule
+1 & -I & \text{elliptic return, memory and past},\\
0  & 0  & \text{parabolic threshold and present},\\
-1 & +I & \text{bidirectional hyperbolic opening and future}.\\
\bottomrule
\end{array}
\]
The TRIT precedes external time: it first classifies the algebraic and
topological type of the transport; an ordered parameterization then publishes it
as a temporal relation. The twofold direction is typed by the exact involution
\begin{equation}
\mathcal C_{\rm rev}(\tau_1,\ldots,\tau_n)
=(-\tau_n,\ldots,-\tau_1),
\qquad \mathcal C_{\rm rev}^2=I,
\label{p12-83-c3-s1:eq:trit-time-reversal}
\end{equation}
which exchanges return and opening and preserves the threshold. On the
hyperbolic sheet, \(P_\pm=(I\pm J_{-1})/2\) separate the two eigendirections, and
reversal exchanges their orientations. Topology and time are thus the two
coordinated readings of the TRIT.

The combinatorial involution \(\mathcal C_{\rm rev}\) and an operatorial
conjugation \(C\) belong, however, to distinct domains. In a
separate dynamic realization, endowed with \(H=H^*\),
\(CJ_EC^{-1}=-J_E\), and \(CHC^{-1}=H\), we fix an internal action scale
\(\hbar_{\rm int}>0\) and write
\(h_{\rm int}:=2\pi\hbar_{\rm int}\). Then
\begin{equation}
U(t)=\exp(-J_EHt/\hbar_{\rm int}),
\qquad CU(t)C^{-1}=U(-t).
\label{p12-83-c3-s1:eq:trit-bidirectional-propagator}
\end{equation}
The intertwiner that transports TRIT words to the dynamical realization fixes
the correspondence between \(\mathcal C_{\rm rev}\) and \(C\). The twofold
local orientation provides the domain of that construction.

\subsection{TPK: category of transport and phase evolution}

The \emph{Topological Phase Kernel} is the transport category over the
observable base of APP that composes digits, words, operators, and prolongation
relations within a single enriched state. Its internal architecture is distributed among
eight operator classes, differentiated by domain, codomain, and function:
\begin{enumerate}[label=\textup{(\roman*)},leftmargin=2.25em]
\item observable support and enriched state;
\item internal sheet selector \(M_{\rm ph}:\{1,\ldots,9\}\to\{+1,-1,0\}\);
\item cardinal translations and visible evolution \(F_{\rm obs}\);
\item readers and quotients nonadic, tritico and of block \(1000\);
\item memory, carry, cylinder, boundary, and genealogical register;
\item periodicity and return
\(C_{12}\hookrightarrow C_{108}\twoheadrightarrow C_9\),
\(\mathcal K_{27}\), and channels \(90/120\);
\item prolongations \(\operatorname{Ext}_r\), cylinders, survival, and
natural bidirectional extension;
\item Archimedean-output, calibrated-exceptional, and dimensional functors.
\end{enumerate}
The return of 27 steps is the composite morphism.
\begin{equation}
\mathcal K_{27}:=F_{\rm obs}^{27},
\qquad \mathcal K_{27}^{4}=F_{\rm obs}^{108}.
\label{p12-83-c3-s1:eq:tpk-kernel-27}
\end{equation}
Its fourth iterate restores the observable phase of the calendar, while the
memory and the complete TPK state continue along their trajectory;
\(\mathcal K_{27}\) is the internal return morphism of \(C_{108}\).
The calendar \(C_{108}\), the nonadic connection, the bilateral extension, and
the readers occupy typed places within this architecture. Let
\(\mathsf P_{\rm APP}\) be the free category
of the directed graph of APP, and let
\(\mathsf P_{\rm APP}^{(2)}=\mathsf P_{\rm APP}\times
\mathsf P_{\rm APP}\) be its simultaneous sum--product reading. Let, moreover,
\[
\mathcal Q_{\rm obs}
=X_9^2\times\{N,E,S,O\}^2\times\Theta_{108},
\qquad \Theta_{108}=\mathbb Z/108\mathbb Z,
\]
where, for the representative
\(\widetilde\theta\in\{0,\ldots,107\}\),
\[
r(\theta)=1+(\widetilde\theta\bmod9),
\qquad g_0(\theta)=[\widetilde\theta]_9,
\]
\[
\tau(\theta)=
\begin{cases}
+1,&r(\theta)\in\{1,4,7\},\\
-1,&r(\theta)\in\{2,5,8\},\\
0,&r(\theta)\in\{3,6,9\},
\end{cases}
\qquad
\mu(\theta)=
\begin{cases}
\Sigma,&\tau(\theta)=+1,\\
\Pi,&\tau(\theta)=-1,\\
\varnothing,&\tau(\theta)=0.
\end{cases}
\]
The calendar label \(\mu=\varnothing\) marks the threshold, and the dynamic memory \(m_k\) of \eqref{p12-83-c3-s1:eq:trit-zero-memory} retains the last active sheet. A configuration \(q=(p^\Sigma,p^\Pi,d^\Sigma,d^\Pi,\theta)\) preserves the two positions, the two cardinal directions, and the calendar index. The visible evolution \(F_{\rm obs}\) activates the sheet indicated by TRIT, transports the corresponding position, and advances \(\theta\). Let \(\mathsf P_{\rm obs}\) denote the free category generated by the arrows \(q\to F_{\rm obs}(q)\). Each arrow induces a pair of APP paths and defines the base functor
\begin{equation}
B:\mathsf P_{\rm obs}\longrightarrow\mathsf P_{\rm APP}^{(2)}.
\label{p12-83-c3-s1:eq:tpk-bivariant-base-functor}
\end{equation}

Over each \(q\) lies the typed fibre
\begin{equation}
\mathcal F_q=
\mathsf O_q\times\mathsf H_q\times\mathsf C_q\times\mathsf S_q
\times\mathsf B_q\times\mathsf\Lambda_q,
\label{p12-83-c3-s1:eq:tpk-fiber}
\end{equation}
Its six factors conserve, respectively, orientation and sheet; residue and
quotient; carry; signature; cylinder and boundary; and ordered register of
transport. Memory and survival are properties of the composition and of
the prolongation relations.

For a route \(r:q\to q'\), the relation
\begin{equation}
\operatorname{Ext}_r\subseteq\mathcal F_q\times\mathcal F_{q'}
\label{p12-83-c3-s1:eq:tpk-extension-relation}
\end{equation}
contains all admissible continuations and satisfies
\[
\Delta_q\subseteq\operatorname{Ext}_{\operatorname{id}_q},
\qquad
\operatorname{Ext}_{r_2}\circ\operatorname{Ext}_{r_1}
\subseteq\operatorname{Ext}_{r_2\circ r_1}.
\]
Transports respect identities and composition:
\begin{equation}
\begin{gathered}
\mathsf H(\operatorname{id}_q)=I,
\quad \mathsf C(\operatorname{id}_q)=I,
\quad \mathsf\Lambda(\operatorname{id}_q)=I,
\quad \kappa(\operatorname{id}_q)=0,\\
\mathsf H(r_2\!\circ r_1)=\mathsf H(r_2)\mathsf H(r_1),
\quad
\mathsf C(r_2\!\circ r_1)=\mathsf C(r_2)\mathsf C(r_1),
\quad
\mathsf\Lambda(r_2\!\circ r_1)=
\mathsf\Lambda(r_2)\mathsf\Lambda(r_1).
\end{gathered}
\label{p12-83-c3-s1:eq:tpk-functorial-transport}
\end{equation}
Along with these functorial laws, the transported coordinates obey
\begin{equation}
\begin{aligned}
h'&=\mathsf H(r)(h),\\
c'&=\mathsf C(r)(c)+\kappa(r),\\
\lambda'&=\mathsf\Lambda(r)(\lambda),\\
\sigma'&=\operatorname{Sig}_{q'}(h',c',\lambda'),
\end{aligned}
\qquad
\kappa(r_2\circ r_1)
=\kappa(r_2)+\mathsf C(r_2)\kappa(r_1).
\label{p12-83-c3-s1:eq:tpk-typed-transport}
\end{equation}
The last identity is the carry cocycle that makes the composition with memory associative.

The TPK is then the category \(\mathsf{TPK}\) whose objects are
\begin{equation}
x=(q,o,h,c,\sigma,C,b,\lambda),
\qquad (o,h,c,\sigma,(C,b),\lambda)\in\mathcal F_q,
\label{p12-83-c3-s1:eq:tpk-enriched-object}
\end{equation}
and whose morphisms \(x\to x'\) are the triples \((r,x,x')\) such that
the coordinates satisfy \eqref{p12-83-c3-s1:eq:tpk-typed-transport} and
\((f_x,f_{x'})\in\operatorname{Ext}_r\), where
\(f_x\in\mathcal F_q\) and \(f_{x'}\in\mathcal F_{q'}\). Composition is
\begin{equation}
(r_2,x',x'')\circ(r_1,x,x')
=(r_2\circ r_1,x,x''),
\label{p12-83-c3-s1:eq:tpk-morphism-composition}
\end{equation}
\Needspace{7\baselineskip}
the identity is \((\operatorname{id}_q,x,x)\), and the functor
\begin{equation}
p:\mathsf{TPK}\longrightarrow\mathsf P_{\rm obs},
\qquad p(x)=q,
\qquad p(r,x,x')=r,
\label{p12-83-c3-s1:eq:tpk-base-functor}
\end{equation}
publishes the visible configuration without identifying states of the same fiber.
The definition fixes a category over a base; it does not presuppose that it is a
categorical fibration nor that every arrow admits a Cartesian lift.
The local projections remain separated by type:
\[
\pi_{\rm TRIT}:\operatorname{Ob}(\mathsf{TPK})
\longrightarrow\mathrm{TRIT}_{\rm vis},
\qquad
\pi_{\rm TRIT}(x)=\tau(\theta),
\]
\[
\pi_{\rm loc}(x)=
\bigl(\pi_{\rm TRIT}(x),\mu(\theta),o,h,c,
\sigma,\lambda,g_0(\theta)\bigr).
\]
The first recovers only the visible TRIT signature. The second further preserves
the calendrical sheet label, orientation, the residue--quotient pair,
carry, transport signature, genealogical register, and phase,
but still forgets the cylinder, its boundary label, and other coordinates of
\(q\). This is why neither the visible TRIT nor its local register
replaces the enriched state of the TPK.

The name encompasses three inseparable functions. \emph{Topological} refers to
the category of paths, incidence, boundary, and prolongation by
cylinders; \emph{Phase} refers to the nonadic index, the oriented calendar, and
holonomy; \emph{Kernel} names the persistent transport kernel that
remains before the state is contracted by a reader. Formally, the
object is the preceding category: an additional topology on objects or
morphisms, or its identification with the linear kernel of a map, requires
independent data.

Consequently, two routes with the same visible endpoint may remain
objects distinct by sheet, quotient, carry, boundary, cylinder, phase, or
register. The TPK does not choose a scalar output: it conserves the common antecedent
from which the Archimedean, exceptional, and dimensional publications depart.

% Estas consecuencias se publican después de la construcción correlativa y
% del paso al límite de los capítulos 4 y 5. Se definen aquí para preservar
% íntegramente su procedencia, pero no se ejecutan antes de que exista el
% continuo al que se aplican.
\long\def\HMTDeferredContinuumConsequences{%
\section{Consequences and publications of the constructed continuum}

\subsection{Terminal Assembly Theorem of the Correlative Constructions Generated by the Common State}

Each depth \(k\) possesses an enriched state \(\mathcal E_k\) with
all its surviving emissions and a single correlative constructor
\begin{equation}
\mathcal O_k:\mathcal E_k\longrightarrow\mathcal Z_k,
\qquad
\widehat x_k\longmapsto
\mathcal O\bigl(\{\mathfrak e_\alpha:
\alpha\in\operatorname{Ch}_k(\widehat x_k)\}\bigr).
\label{p12-83-c3-s1:eq:five-features-single-constructor}
\end{equation}
The value of \(\mathcal O_k\) simultaneously preserves the spectral,
solenoidal, gauge, coherent, and determinantal coordinates,
\begin{equation}
(\mathrm{sp},\mathrm{sol},\mathrm{gau},\mathrm{coh},\mathrm{det}),
\label{p12-83-c3-s1:eq:five-features-coordinates}
\end{equation}
as five characteristics of the same state and of the same genealogical register. The
truncations
\[
u^{\mathfrak e}_{\ell,k}:\mathcal E_\ell\to\mathcal E_k,
\qquad
u^{\mathcal O}_{\ell,k}:\mathcal Z_\ell\to\mathcal Z_k
\]
satisfy the unique naturality
\begin{equation}
u^{\mathcal O}_{\ell,k}\circ\mathcal O_\ell
=\mathcal O_k\circ u^{\mathfrak e}_{\ell,k},
\qquad \ell\ge k.
\label{p12-83-c3-s1:eq:five-features-naturality}
\end{equation}
Therefore, the corresponding action passes to the limit as
\begin{equation}
\mathcal O_\infty:
\mathcal C_{\rm cont}^{\rm disc}:=\varprojlim_k\mathcal E_k
\longrightarrow
\mathcal Z_\infty:=\varprojlim_k\mathcal Z_k.
\label{p12-83-c3-s1:eq:five-features-infinite-operator}
\end{equation}
The five terminal applications are seen as intrinsic manifestations of that single value:
\begin{equation}
\boxed{
\mathfrak G_T^{\rm int}:=
\operatorname{View}_T\!\left(
\mathcal O_\infty(\mathcal C_{\rm cont}^{\rm disc})\right),
\quad
T\in\{\mathrm{sp},\mathrm{sol},\mathrm{gau},
\mathrm{coh},\mathrm{det}\}.}
\label{p12-83-c3-s1:eq:five-features-views}
\end{equation}
Simultaneity precedes the specialized resolutions: each view exposes
the required coordinate and the complete state conserves the other four.

\subsection{Survival and inverse limit}

Let \(\operatorname{Surv}^{(\chi)}_N\) be the set of histories of
depth \(N\) that simultaneously satisfy the active laws of channel
\(\chi\) and admit prolongation. We define the compatible truncations
\[
\tau_{N+1,N}:\operatorname{Surv}^{(\chi)}_{N+1}
\longrightarrow\operatorname{Surv}^{(\chi)}_N
\]
These preserve the enriched prefix. The genealogical continuum is the limit.
\begin{equation}
X_\chi=\varprojlim_N
\bigl(\operatorname{Surv}^{(\chi)}_N,\tau_{N+1,N}\bigr),
\qquad
\tau_{N+1,N}(x_{N+1})=x_N.
\label{p12-83-c3-s1:eq:inverse-survival}
\end{equation}
The global space of surviving sections and one of its enriched sections
are denoted, respectively, by
\begin{equation}
X_\infty:=\bigsqcup_\chi X_\chi,
\qquad
\widehat\Omega_\infty\in X_\infty.
\label{p12-83-c3-s1:eq:global-survival-space}
\end{equation}
The global Archimedean reader is
\begin{equation}
P_{\rm arq}:=\operatorname{ev}_{\mathbb R}:
X_\infty\longrightarrow\mathbb R;
\label{p12-83-c3-s1:eq:archimedean-reader}
\end{equation}
its restriction to each channel evaluates the unitary intersection of the compatible cylinders of the corresponding section. When the exceptional reader intervenes, \(X_\infty^{\rm cal}\subseteq
X_\infty\) denotes the subspace in which the calibration type of \eqref{p12-83-c3-s1:eq:exceptional-calibration} has been fixed. Survival determines the section by internal compatibility: every history must be prolongable in the complete system. The numerical reader intervenes afterward to certify and publish the intersection of the already compatible descendant.

\subsection{Joint nonadic connection, holonomy, and vertical memory}

At depth \(K\), the holonomy of the nonadic return is initially the typed correspondence
\begin{equation}
\Gamma_{9,K}=\operatorname{Hol}_{C_9}(\nabla_{\rm TPK}):
\widehat{\mathcal X}_K\rightrightarrows
\widehat{\mathcal X}_{K+9}.
\label{p12-83-c3-s1:eq:gamma9}
\end{equation}
Its iterates are relational compositions. Only when each state has
a unique compatible prolongation is it restricted to a map, and only when that
map is invertible does it constitute a monodromy. This distinction conserves all
surviving extensions before selecting a section.
After nine phases the local phase returns, while the complete state advances:
\[
g(K+9)=g(K),\qquad q_{{\rm mem},K+9}=q_{{\rm mem},K}+1,
\qquad \widehat x_{K+9}\ne\widehat x_K\ \text{in general}.
\]
Block, cylinder, prefix, carry, and memory continue to evolve. This
is the difference between visible periodicity and return with biography.

The transported state at level (K) can be exhibited as
\begin{equation}
\widehat x_K=
(B_K,C_K,p_K,c_K,\Sigma_K,W_K,g(K),q_{{\rm mem},K},s_K),
\label{p12-83-c3-s1:eq:enriched-state-k}
\end{equation}
where each coordinate has its own transport law. The phase return is represented on the catalog.
\[
\mathcal U_6^{\rm cat}\cong\mathfrak S_+\times\mathfrak S_\times,
\qquad |\mathfrak S_+|=18,
\qquad |\mathfrak S_\times|=26.
\]
Conditional averages
\[
(E_+f)(s,e)=\frac1{18}\sum_{s'\in\mathfrak S_+}f(s',e),
\qquad
(E_\times f)(s,e)=\frac1{26}\sum_{e'\in\mathfrak S_\times}f(s,e')
\]
commute. The phase word
\(\tau=(+1,-1,0,+1,-1,0,+1,-1,0)\) applies
\(E_+,E_\times,I\) in that order, and one revolution closes through
\begin{equation}
\mathcal K_{\rm ph}^{\,9}
=(E_+E_\times)\otimes I_{\ell^2(C_9)},
\qquad
\bigl\langle\delta_u,E_+E_\times\delta_v\bigr\rangle
=\frac1{468}\quad(u,v\in\mathcal U_6^{\rm cat}).
\label{p12-83-c3-s1:eq:joint-nine-holonomy}
\end{equation}
The sum and product sheets do not act as independent filters: they are two
components of a single holonomy that closes only after traversing the
complete cycle.

The compression--expression colligation separates what returns from what remains
as memory. If (J) embeds the history space into the phase
space, (T) is the surviving compression and \(\eta\) the defect channel,
then
\begin{equation}
\boxed{
U_{\Gamma_9}J=JT+\eta,
\qquad J^*\eta=0,
\qquad I-T^*T=\eta^*\eta.}
\label{p12-83-c3-s1:eq:compression-expression}
\end{equation}
The first equality decomposes transport into compressed return and memory
expression; the second proves their orthogonality; the third measures exactly
the information that cannot fit within the periodic shadow. The central
statement is thereby formalized: the phase returns, not the complete state.

The vertical memory is organized through the cyclic extension oriented
\begin{equation}
0\longrightarrow C_{12}\xrightarrow{\,\iota\,}C_{108}
\xrightarrow{\,p\,}C_9\longrightarrow0,
\qquad
c(r,r')=\left[\left\lfloor\frac{r+r'}9\right\rfloor\right]_{12}.
\label{p12-83-c3-s1:eq:c108-extension}
\end{equation}
Here \(\iota([b]_{12})=[9b]_{108}\) and
\(p([\theta]_{108})=[\theta]_9\). The extension does not split:
\(C_{108}\) has exponent 108, whereas
\(C_{12}\times C_9\) has exponent 36.
The cocycle records the carry from each turn of the base \(C_9\) in the
dodecaphase memory \(C_{12}\). \(C_{108}\) is therefore an oriented
phase--memory calendar. Only after the dimensional chart is that
extension realized as a finite 108-step quantum clock.

The local MD lift carrying the TRIT signature is the
topological--temporal unit of that genealogical crystal.
On
\(\mathcal H_{\rm clk}=\mathbb C^{108}\), with basis
\(\{|j\rangle:j\in\mathbb Z/108\mathbb Z\}\), let
\[
\zeta=e^{2\pi\ii/108},\qquad
|\widetilde k\rangle=\frac1{\sqrt{108}}
\sum_{j=0}^{107}\zeta^{jk}|j\rangle,
\qquad U_{\rm clk}|j\rangle=|j+1\rangle.
\]
\Needspace{12\baselineskip}
With \(\omega_0=2\pi/(108t_0)\), the Hamiltonian
\begin{equation}
H_{\rm clk}=\hbar_{\rm int}\omega_0
\sum_{k=0}^{107}k\,|\widetilde k\rangle\langle\widetilde k|
\label{p12-83-c3-s1:eq:trit-time-crystal-hamiltonian}
\end{equation}
generates exactly
\begin{equation}
U_{\rm step}=\exp\!\left(
-\frac{\ii H_{\rm clk}t_0}{\hbar_{\rm int}}\right)=U_{\rm clk},
\qquad U_{\rm step}^{108}=I,
\label{p12-83-c3-s1:eq:trit-time-crystal-period}
\end{equation}
If \(Z_{\rm clk}|j\rangle=\zeta^j|j\rangle\), then
\begin{equation}
\langle j_0|U_{\rm step}^{-n}Z_{\rm clk}U_{\rm step}^{n}|j_0\rangle
=\zeta^{j_0+n},
\label{p12-83-c3-s1:eq:trit-time-crystal-observable}
\end{equation}
whose orbit has primitive period 108. In HMT terminology, this is the
internal genealogical time crystal: the observable calendar closes,
\(C_{12}\) records the circuit, and the enriched TPK state need not
return. Its subsequent realization as a macroscopic physical phase also requires
a Floquet protocol, primitive subharmonic response, robustness,
persistence, and a protection mechanism; that chart does not alter the finite
theorem of unitary periodicity.

\subsection{Archimedean publication, genealogical number and measurement}

The Archimedean publication uses semi-open ternary cylinders.
\begin{equation}
I_{N,L}=\left[\frac{N}{3^L},\frac{N+1}{3^L}\right),
\qquad \operatorname{diam}I_{N,L}=3^{-L}\longrightarrow0.
\label{p12-83-c3-s1:eq:ternary-cylinders}
\end{equation}
A compatible section determines nested cylinders and, by completeness, a
unique real point. The HMT number is the section itself,
\begin{equation}
x=(x_N)_N\in X_\chi
=\varprojlim_N\operatorname{Surv}^{(\chi)}_N,
\label{p12-83-c3-s1:eq:hmt-number-as-section}
\end{equation}
not the tuple formed by it and its evaluations. The value appears when a
reader forgets part of the genealogy.
Therefore, precision is depth: measuring means descending to a
finer cylinder while the causal prefix remains in the antecedent.
The real line is reconstructed as the Archimedean quotient of a space of
much richer histories.

\subsection{Dimensional realization of the HMT state: screen map, observable magnitude, and spectral chart}

Dimensional Mechanics begins where the genealogical state couples to a
realization boundary. With a field \(K_n\) fixed, let \(M_n(x)\) be the module
generated by the prolongation channels of a state \(x\), let
\(\mathcal R_n(x)\) be the submodule of its relations, and let \(\beta_n(x)\) be the
conserved boundary datum. The dimensional screen is
\begin{equation}
\Sigma_n(x)=\bigl(K_n,M_n(x),\beta_n(x),\mathcal R_n(x)\bigr),
\qquad
d_{\Sigma_n}(x)=\dim_{K_n}\frac{M_n(x)}{\mathcal R_n(x)}.
\label{p12-83-c3-s1:eq:md-screen}
\end{equation}
If \(M_n(x)\simeq K_n^{g_n}\) and the columns of \(R_n(x)\) generate the
relations, then
\begin{equation}
\boxed{d_{\Sigma_n}(x)=g_n-\operatorname{rank}_{K_n}R_n(x).}
\label{p12-83-c3-s1:eq:md-rank}
\end{equation}
Dimension is not attached as a label: it is the number of independent
channels that survive after imposing the boundary relations.

A physical quantity is a section of a dimensional line \(L_D\); a unit is a chart \(\chi_u:L_D\to\mathbb R\) that publishes its coordinate. Hence
\begin{equation}
\text{TPK state}\longrightarrow\Sigma_n(x)
\longrightarrow s_D\in\Gamma(L_D)
\xrightarrow{\;\chi_u\;}[s_D]_u\in\mathbb R
\label{p12-83-c3-s1:eq:md-measurement-chain}
\end{equation}
separates structure, dimensional type, and measured number exactly. Measuring
consists in coupling a section to a reversible apparatus, imposing
edge compatibility, and applying the reader; changing units changes
\(\chi_u\), not the section or its genealogy.

\subsection{Double projection mathematical and dimensional realization}

The exceptional branch preserves gauge data that the real publication does not
need. Let \(\mathscr C_{\rm exc}\) be the space of gauges, and let
\begin{equation}
\mathfrak E:
X_\infty^{\rm cal}\times\mathscr C_{\rm exc}
\longrightarrow\mathbf{Exc}
\label{p12-83-c3-s1:eq:exceptional-calibrated-reader}
\end{equation}
the incidence reader. We fix once
\begin{equation}
\begin{aligned}
\chi_{\rm exc}=\bigl(&\text{projective order},\ \text{Paley gauge},
 \ \text{dodecad},\\
 &\text{incidence alignment},\ \text{lattice chamber}\bigr)
 \in\mathscr C_{\rm exc}
\end{aligned}
\label{p12-83-c3-s1:eq:exceptional-calibration}
\end{equation}
and abreviamos
\(\mathfrak E_{\chi_{\rm exc}}(x):=\mathfrak E(x,\chi_{\rm exc})\).
The gauge initial it transports by the history; not vuelve to elegirse in
each depth.

The same calibrated limit then admits two readers with distinct codomains:
\begin{equation}
\mathbb R
\xleftarrow{\;\operatorname{ev}_{\mathbb R}\;}
X_\infty^{\rm cal}
\xrightarrow{\;\mathfrak E_{\chi_{\rm exc}}\;}
\mathbf{Exc}.
\label{p12-83-c3-s1:eq:double-projection}
\end{equation}
The first branch contracts cylinders to a continuous value. The second
preserves incidences, orientations, gauge, and boundary. Both descend in
parallel from the same source state. This bifurcation jointly explains
metric continuity and the emergence of exceptional symmetries. The second
branch is a calibrated reconstruction: it is exact once
\(\chi_{\rm exc}\) has been fixed, but the bare state does not retrospectively select
the five data composing that gauge.

The non-serializing dimensional realization does not serialize those two branches. It is a third reader of the same antecedent:
\begin{equation}
P_{\rm dim}:X_\infty\longrightarrow
\mathcal R_{\rm MD},
\qquad
\mathcal R_{\rm MD}=
\{\text{action, mass, field, geometry, information}\}.
\label{p12-83-c3-s1:eq:dimensional-publication}
\end{equation}
In this manner, a number, a symmetry, and a physical magnitude may share a
genealogy without being confused with or reduced to one another.

The thesis of the continuum is thus formulated precisely: Archimedean
reality is the zero-radius shadow of infinitely refinable discrete histories;
what the shadow omits reappears as memory, incidence,
geometry, and physical structure.

% RESULTADO_RECUPERADO. Owner íntegro de la dinámica del cociente, el retorno
% observable y los cociclos de modo y orientación. Se subordina a la única
% construcción APP--TRIT--TPK sin abrir un origen alternativo.
\begin{HMTScientificOwnerSubordinated}
\input{sections/hmt/05b_extensiones_dinamica_cociente}
\end{HMTScientificOwnerSubordinated}
}
